% GENERATED FILE. Edit canonical lessons, then run npm run generate:books.
% canonical-source-hash: fnv1a64:0cb05c07b1a1edeb
% canonical-lessons: GE-C508-siegen, GE-C508-sieger, GE-C508-verlierer, GE-C508-gegner, GE-C508-turnier

\chapter{Sieger und Verlierer}
\label{ch:ge-sieger-und-verlierer}
\hlchaptermodality{508}

\begin{chapteropening}
\textbf{By the end of this chapter:} \emph{I can talk about winning and losing.}

Complete the last lesson of chapter 508: I can talk about winning and losing.
\end{chapteropening}

\section[siegen]{siegen — to win, to be victorious}

\label{lesson:GE-C508-siegen}



\begin{quote}
Before the new one: say the German for tied game, then the German for offside.
\end{quote}

\subsection*{You'll want to know: siegen}
\textbf{siegen} — \textquotedblleft{}to win, to be victorious\textquotedblright{}.

\begin{quote}
Unsere Mannschaft hat mit drei zu null gesiegt.
\end{quote}

\begin{quote}
\emph{Our team won three-nil.}
\end{quote}

More formal than \emph{gewinnen}. \emph{Der Sieg} — victory.

\begin{cousinweb}
\emph{siegen} is built on \emph{Sieg} \textquotedblleft{}victory\textquotedblright{} (EtymDB), inherited from Middle High German \emph{sige} and Old High German \emph{sigu}, from Proto-Germanic \emph{segaz}.
\end{cousinweb}

\begin{grammarlens}[title={What you've built}]
One more word for telling, explaining and describing at B1.
\end{grammarlens}

\subsection*{Your turn}
\begin{itemize}
\raggedright
  \item \emph{Say it:} \emph{siegen}
  \item \emph{Say it:} \emph{siegen}, once more
  \item \emph{Say it:} a sentence of your own with \emph{siegen}
  \item \emph{Recall:} say the German for tied game, then the German for offside, then say \emph{siegen} again
\end{itemize}

\begin{tcolorbox}[breakable,colback=teal!4,colframe=teal!35!black,title={Before you move on}]
What does siegen mean? (\textquotedblleft{}To win, to be victorious\textquotedblright{}.) Say it once more.
\end{tcolorbox}
\section[der Sieger]{der Sieger — winner}

\label{lesson:GE-C508-sieger}



\begin{quote}
Before the new one: say the German for offside, then the German for to win.
\end{quote}

\subsection*{You'll want to know: der Sieger}
\textbf{der Sieger} — \textquotedblleft{}winner\textquotedblright{}.

\begin{quote}
Am Ende war mein kleiner Bruder der Sieger.
\end{quote}

\begin{quote}
\emph{In the end my little brother was the winner.}
\end{quote}

Plural: \emph{die Sieger}. A woman is \emph{die Siegerin}.

\begin{cousinweb}
\emph{Sieger} is built on \emph{siegen} with \emph{-er}; \emph{Sieg} \textquotedblleft{}victory\textquotedblright{} is inherited from Old High German \emph{sigu}, from Proto-Germanic \emph{segaz} (EtymDB).
\end{cousinweb}

\begin{grammarlens}[title={What you've built}]
One more word for telling, explaining and describing at B1.
\end{grammarlens}

\subsection*{Your turn}
\begin{itemize}
\raggedright
  \item \emph{Say it:} \emph{der Sieger}
  \item \emph{Say it:} \emph{der Sieger}, once more
  \item \emph{Say it:} a sentence of your own with \emph{der Sieger}
  \item \emph{Recall:} say the German for offside, then the German for to win, then say \emph{der Sieger} again
\end{itemize}

\begin{tcolorbox}[breakable,colback=teal!4,colframe=teal!35!black,title={Before you move on}]
What does der Sieger mean? (\textquotedblleft{}Winner\textquotedblright{}.) Say it once more.
\end{tcolorbox}
\section[der Verlierer]{der Verlierer — loser}

\label{lesson:GE-C508-verlierer}



\begin{quote}
Before the new one: say the German for to win, then the German for winner.
\end{quote}

\subsection*{You'll want to know: der Verlierer}
\textbf{der Verlierer} — \textquotedblleft{}loser\textquotedblright{}.

\begin{quote}
Der Verlierer gratulierte dem Sieger.
\end{quote}

\begin{quote}
\emph{The loser congratulated the winner.}
\end{quote}

Plural: \emph{die Verlierer}. \emph{Ein schlechter Verlierer} — a bad loser.

\begin{cousinweb}
\emph{Verlierer} is built on \emph{verlieren} \textquotedblleft{}to lose\textquotedblright{}, inherited from Old High German \emph{firliosan}, from Proto-Germanic \emph{fraleusaną} (EtymDB), with \emph{-er}.
\end{cousinweb}

\begin{grammarlens}[title={What you've built}]
One more word for telling, explaining and describing at B1.
\end{grammarlens}

\subsection*{Your turn}
\begin{itemize}
\raggedright
  \item \emph{Say it:} \emph{der Verlierer}
  \item \emph{Say it:} \emph{der Verlierer}, once more
  \item \emph{Say it:} a sentence of your own with \emph{der Verlierer}
  \item \emph{Recall:} say the German for to win, then the German for winner, then say \emph{der Verlierer} again
\end{itemize}

\begin{tcolorbox}[breakable,colback=teal!4,colframe=teal!35!black,title={Before you move on}]
What does der Verlierer mean? (\textquotedblleft{}Loser\textquotedblright{}.) Say it once more.
\end{tcolorbox}
\section[der Gegner]{der Gegner — opponent}

\label{lesson:GE-C508-gegner}



\begin{quote}
Before the new one: say the German for winner, then the German for loser.
\end{quote}

\subsection*{You'll want to know: der Gegner}
\textbf{der Gegner} — \textquotedblleft{}opponent\textquotedblright{}.

\begin{quote}
Der Gegner war stärker als wir.
\end{quote}

\begin{quote}
\emph{The opponent was stronger than us.}
\end{quote}

Plural: \emph{die Gegner}. Also an opponent of an idea: \emph{ein Gegner der Reform}.

\begin{cousinweb}
\emph{Gegner} comes from Middle Low German \emph{gēgenēre} (EtymDB), built on \emph{gegen} \textquotedblleft{}against\textquotedblright{}.
\end{cousinweb}

\begin{grammarlens}[title={What you've built}]
One more word for telling, explaining and describing at B1.
\end{grammarlens}

\subsection*{Your turn}
\begin{itemize}
\raggedright
  \item \emph{Say it:} \emph{der Gegner}
  \item \emph{Say it:} \emph{der Gegner}, once more
  \item \emph{Say it:} a sentence of your own with \emph{der Gegner}
  \item \emph{Recall:} say the German for winner, then the German for loser, then say \emph{der Gegner} again
\end{itemize}

\begin{tcolorbox}[breakable,colback=teal!4,colframe=teal!35!black,title={Before you move on}]
What does der Gegner mean? (\textquotedblleft{}Opponent\textquotedblright{}.) Say it once more.
\end{tcolorbox}
\section[das Turnier]{das Turnier — tournament}

\label{lesson:GE-C508-turnier}



\begin{quote}
Before the new one: say the German for loser, then the German for opponent.
\end{quote}

\subsection*{You'll want to know: das Turnier}
\textbf{das Turnier} — \textquotedblleft{}tournament\textquotedblright{}.

\begin{quote}
Bei dem Turnier spielen zwölf Mannschaften.
\end{quote}

\begin{quote}
\emph{Twelve teams are playing in the tournament.}
\end{quote}

Plural: \emph{die Turniere}.

\begin{cousinweb}
\emph{Turnier} goes back to Middle High German \emph{turnier}, usually traced to Old French \emph{tornei} \textquotedblleft{}a joust\textquotedblright{}, built on \emph{torner} \textquotedblleft{}to turn\textquotedblright{}, from Latin \emph{tornāre}; EtymDB lists the German word without a source.
\end{cousinweb}

\begin{grammarlens}[title={What you've built}]
One more word for telling, explaining and describing at B1.
\end{grammarlens}

\subsection*{Your turn}
\begin{itemize}
\raggedright
  \item \emph{Say it:} \emph{das Turnier}
  \item \emph{Say it:} \emph{das Turnier}, once more
  \item \emph{Say it:} a sentence of your own with \emph{das Turnier}
  \item \emph{Recall:} say the German for loser, then the German for opponent, then say \emph{das Turnier} again
\end{itemize}

\begin{tcolorbox}[breakable,colback=teal!4,colframe=teal!35!black,title={Before you move on}]
What does das Turnier mean? (\textquotedblleft{}Tournament\textquotedblright{}.) Say it once more.
\end{tcolorbox}
